\section{Evaluation Metrics}
\label{sec:performance_metrics}

In this study, we employed a comprehensive set of evaluation metrics to assess the performance of our models on the \named judgment prediction and PredEx explanation test datasets. We report Macro Precision, Macro Recall, Macro F1, and Accuracy for judgment prediction, and we use both quantitative and qualitative methods to evaluate the quality of explanations generated by the model.

\begin{enumerate}
    \item \textbf{Lexical-based Evaluation:} We utilized standard lexical similarity metrics, including Rouge scores (Rouge-1, Rouge-2, and Rouge-L) \cite{lin-2004-rouge}, BLEU \cite{papineni-etal-2002-bleu}, and METEOR \cite{banerjee-lavie-2005-meteor}. These metrics measure the overlap and order of words between the generated explanations and the reference texts, providing a quantitative assessment of the lexical accuracy of the model outputs.

    \item \textbf{Semantic Similarity-based Evaluation:} To capture the semantic quality of the generated explanations, we employed BERTScore \cite{BERTScore}, which measures the semantic similarity between the generated text and the reference explanations. Additionally, we used BLANC \cite{blanc}, a metric that estimates the quality of generated text without a gold standard, to evaluate the model's ability to produce semantically meaningful and contextually relevant explanations.
    
    \item \textbf{Expert Evaluation:} Human evaluation was a critical component of our assessment framework. Legal experts reviewed the explanations generated by the models, rating them on a 1–5 Likert scale based on criteria such as accuracy, relevance, and completeness. A rating of 1 indicates that the information is irrelevant, while a rating of 5 signifies that the explanation is superior to the expert's own explanation. The full description of the rating scores can be found in Appendix \ref{section:rating-score-description} - Rating Score Description, which is adapted from \cite{nigam2024legaljudgmentreimaginedpredex}.

\end{enumerate}
